\subsolution{Интерференция звука}{3.0}

Введем декартову систему координат. Пусть линия пересечения двух стен совпадает с осью $z$, а сами стены лежат в плоскостях $x=0$ и $y=0$. Волна будет в квадранте $x \ge 0$, $y \ge 0$. Звуковые волны продольные, поэтому молекулы будут колебаться вдоль направления волнового вектора. Падающая плоская волна движется под углом $45^\circ$ к обеим стенам по направлению к началу координат. Ее волновой вектор $\vec{k}_1$ имеет компоненты:
$$
k_x = k_y = k \cos(45^\circ) = \frac{k}{\sqrt{2}} = \frac{\sqrt{2}\pi}{\lambda}
$$
Дальше будем обозначать \(k_0 = \pi \sqrt{2} / \lambda\). Тогда $\vec{k}_1 = (-k_0, -k_0)$. Вектор скорости падающей волны:

\begin{eq}[points=0.4]
    \vec{v}_1(x,y,t) = \left(-\frac{v}{\sqrt{2}}, -\frac{v}{\sqrt{2}}\right) \cos\left(\omega t + \frac{\sqrt{2}\pi}{\lambda} x + \frac{\sqrt{2}\pi}{\lambda} y\right),
\end{eq}

\criterion[type=remark]{\textbf{0.2 + 0.2 } за верную амплитуду вектора и зависимость фазы от координат}

При отражении звуковых волн фаза не меняется. Тогда в пространстве будет суперпозиция 3 разных отраженных волн:

    \textbf{Волна, отраженная от стены $x=0$:} $\vec{k}_2 = (k_0, -k_0)$. Знак $x$-компоненты скорости меняется на противоположный.

    \begin{eq}[points=0.1]
        \vec{v}_2(x,y,t) = \left(\frac{v}{\sqrt{2}}, -\frac{v}{\sqrt{2}}\right) \cos(\omega t - k_0 x + k_0 y),
    \end{eq}

    \textbf{Волна, отраженная от стены $y=0$:} $\vec{k}_3 = (-k_0, k_0)$. Знак $y$-компоненты скорости меняется на противоположный.

    \begin{eq}[points=0.1]
        \vec{v}_3(x,y,t) = \left(-\frac{v}{\sqrt{2}}, \frac{v}{\sqrt{2}}\right) \cos(\omega t + k_0 x - k_0 y),
    \end{eq}

    \textbf{Волна, отраженная от обеих стен (от угла):} $\vec{k}_4 = (k_0, k_0)$.

    \begin{eq}[points=0.1]
        \vec{v}_4(x,y,t) = \left(\frac{v}{\sqrt{2}}, \frac{v}{\sqrt{2}}\right) \cos(\omega t - k_0 x - k_0 y),
    \end{eq}

    \criterion[type=remark]{Выше баллы \textbf{только} за верную амплитуду вектора \textbf{и} фазу}

Найдем суммарную скорость молекул $\vec{v} = \vec{v}_1 + \vec{v}_2 + \vec{v}_3 + \vec{v}_4$, сложив компоненты отдельно.
Для $x$-компоненты:
$$
v_x = \frac{v}{\sqrt{2}} \big[ -\cos(\omega t + k_0 x + k_0 y) + \cos(\omega t - k_0 x + k_0 y) - \cos(\omega t + k_0 x - k_0 y) + \cos(\omega t - k_0 x - k_0 y) \big]
$$

Сгруппируем слагаемые с одинаковой фазой по $y$:
$$
v_x = \frac{v}{\sqrt{2}} \big[ \big(\cos(\omega t + k_0 y - k_0 x) - \cos(\omega t + k_0 y + k_0 x)\big) + \big(\cos(\omega t - k_0 y - k_0 x) - \cos(\omega t - k_0 y + k_0 x)\big) \big]
$$

Используя формулу разности косинусов $\cos(A-B) - \cos(A+B) = 2\sin A \sin B$:
$$
v_x = \frac{v}{\sqrt{2}} \big[ 2\sin(\omega t + k_0 y)\sin(k_0 x) + 2\sin(\omega t - k_0 y)\sin(k_0 x) \big]
$$

Вынесем общий множитель и применим формулу суммы синусов $\sin(A+B) + \sin(A-B) = 2\sin A \cos B$:
$$
v_x = \sqrt{2}v \sin(k_0 x) \big[ \sin(\omega t + k_0 y) + \sin(\omega t - k_0 y) \big] = 2\sqrt{2}v \sin(k_0 x) \cos(k_0 y) \sin(\omega t)
$$

В силу симметрии задачи (при замене $x \leftrightarrow y$), $y$-компонента скорости имеет вид:
\begin{eq}[points=0.8, tail={, или эквивалентное уравнение двойной стоячей волны для хотя бы одной из компонент скорости}]
    v_y = 2\sqrt{2}v \cos(k_0 x) \sin(k_0 y) \sin(\omega t),
\end{eq}

Как мы видим, компоненты скорости имеют вид "двойной" стоячей волны.

Квадрат модуля скорости молекул $|\vec{v}|^2 = v_x^2 + v_y^2$:

\begin{align*}
    |\vec{v}|^2 & = 8v^2 \sin^2(\omega t) \big[ \sin^2(k_0 x) \cos^2(k_0 y) + \cos^2(k_0 x) \sin^2(k_0 y) \big] = 8v^2 \sin^2(\omega t) f(x,y)
\end{align*}

Взяв производную от выражения в квадратной скобке по двум разным координатам:

\begin{align*}
    \frac{\partial f(x,y)}{\partial x} & =  2k_0 \sin(k_0x) \cos(k_0 x) \cos^2(k_0 y) - 2 k_0 \cos(k_0 x ) \sin(k_0 x) \sin^2(k_0y)
    \\
    \frac{\partial f(x,y)}{\partial x} & = k_0 \sin(2 k_0 x) \cos(2 k_0y) = 0
\end{align*}

Точно такое же условие и для производной по \(y\). Получаем условия для экстремума:

\begin{gather*}
    \sin(2 k_0 x) \cos(2 k_0y) = 0, \quad  \sin(2 k_0 y) \cos(2 k_0x) = 0
    \\
     \left( x = \frac{\lambda}{2 \sqrt{2}} n \ \ \text{или} \ \  y =  \frac{\lambda}{2 \sqrt{2}} \left(\frac{1}{2} + m \right)\right) \ \ \text{и} \ \ \left( y = \frac{\lambda}{2 \sqrt{2}} n \ \ \text{или} \ \  x =  \frac{\lambda}{2 \sqrt{2}} \left(\frac{1}{2} + m \right)\right)
\end{gather*}

Из условий экстремума понятно что максимум \(f(x,y)\) равен \(1\) (это можно было и показать по другому). Тогда сразу же можно найти максимальную скорость:

\begin{eq}[points=0.5]
    v_{\max} = \sqrt{8v^2} = 2\sqrt{2}v,
\end{eq}

Координаты точек максимума можно подобрать:
\begin{eq}[points=1.0]
    \left( \frac{\lambda}{2 \sqrt{2}} \cdot n, \frac{\lambda}{2 \sqrt{2}} \cdot m \right) \ \ \text{где ровно  одно из \(n,m\) нечётное}.
\end{eq}

\criterion[type=remark]{\textbf{1.0} за полностью верный и обоснованный ответ для координат максимумов}

\criterion[type=remark]{\textbf{0.4} за частично верный ответ, где видно что точки максимума будут на прямоугольной \textbf{бесконечной сетке}}
